\chapter{Разработка программного обеспечения} \label{ch3}

% не рекомендуется использовать отдельную section <<введение>> после лета 2020 года
%\section{Введение} \label{ch3:intro}

В данной главе перечислены инструменты, применявшиеся в процессе решения задачи.
\section{Предобработка данных}
Предобработка данных производилась средствами языка Python 3.11.8.
Применены библиотеки:
\begin{itemize}
	\item \textbf{pandas} - для форматирования, валидации и препроцессинга объемов данных.
	\item \textbf{csv} - для работы с исходными файлами, содержащими данные, в формате csv.
\end{itemize}
\section{Алгоритмы оптимизации}
\begin{itemize}
	\item \textbf{pyGAD}, \textbf{pyswarms}, \textbf{cma} - библиотеки, содержащие абстракции для построения ГА, МРЧ и CMA-ES соответственно.
\end{itemize}
\section{Симуляция движения}
Реализована на базе фреймворка SUMO.
\begin{itemize}
	\item \textbf{dfrouter} - скрипт для создания путей (маршрутов) по информации о местоположении детекторов, с которых были собраны данные о трафике.
	\item \text{osmWebWizard} - инструмент для "one-click" конвертации реальной дорожной карты в $G(V, E)$.
	\item \textbf{netedit} - инструмент для редактирования результатов работы osmWebWizard, а также добавления светофорной логики, мест установки детекторов на карту.
	
\end{itemize}
\section{Работа с графическим представлением задачи} \label{ch3:sec1}
Для графического представления задачи был использован инструмент
\section{Название параграфа} \label{ch3:sec2}

%\FloatBarrier % заставить рисунки и другие подвижные (float) элементы остановиться


\section{Выводы} \label{ch3:conclusion}

Текст выводов по главе \thechapter.


%% Вспомогательные команды - Additional commands
%
%\newpage % принудительное начало с новой страницы, использовать только в конце раздела
%\clearpage % осуществляется пакетом <<placeins>> в пределах секций
%\newpage\leavevmode\thispagestyle{empty}\newpage % 100 % начало новой страницы